% ShiftLens competition report (max 5 A4 pages, pdflatex)
\documentclass[11pt,a4paper]{article}
\usepackage[margin=2.2cm]{geometry}
\usepackage[T1]{fontenc}
\usepackage{lmodern}
\usepackage{microtype}
\usepackage{amsmath}
\usepackage{booktabs}
\usepackage{xcolor}
\usepackage{enumitem}
\usepackage[colorlinks=true,linkcolor=black,urlcolor=blue!45!black]{hyperref}

\definecolor{slblue}{RGB}{20,60,120}
\setlength{\parindent}{0pt}
\setlength{\parskip}{4pt plus 1pt}
\setlist{nosep,leftmargin=1.5em}

\title{\textcolor{slblue}{\textbf{ShiftLens}}\\[2pt]
\large Agentic Workforce Analytics for Productivity and Wellbeing}
\author{ShiftLens Team}
\date{October 2026}

\begin{document}
\maketitle

\begin{abstract}
\noindent
ShiftLens is an agentic workforce-analytics system for shift teams on a
production floor. Every day it ingests each worker's wearable and
environmental sensor streams, micro-survey answers, and camera-based
object-detection events, and writes an agent-authored daily report per
worker; at month's end it correlates wellbeing and working conditions with
productivity and serves an offline-capable dashboard of findings. All
metrics, flags, correlations, and findings are produced by a deterministic
Python core --- the guiding principle is \emph{the LLM phrases; it never
decides}: an optional local LLM only rewrites the prose summary, and
everything works with the LLM offline. A seeded simulator with planted
causal structure stands in for real feeds in the demo and provides the
ground truth used to validate the analytics end-to-end.
\end{abstract}

\section{Introduction and problem}
Supervisors of shift teams --- assemblers, packers, inspectors, forklift
operators --- sit on fragmented data. Wearables report heart rate, steps,
and station time; environmental sensors log noise, temperature, and air
quality; camera systems emit object-detection events such as PPE
non-compliance, phone use, idle periods, and safety-zone violations; short
daily surveys capture fatigue and mood. Each stream lives in its own silo,
and none of it is joined up with the numbers that matter: units produced
and defects.

Three pain points follow. \textbf{Fragmentation}: no single place shows one
worker's day across all sources. \textbf{Latency}: productivity is typically
reviewed in aggregate at month's end, so problems --- a noisy line, a heat
wave, a fatigued crew --- are discovered weeks after they start costing
output and quality. \textbf{No wellbeing--output link}: fatigue and mood
are collected but never connected to output, so the cost of poor conditions
stays invisible.

ShiftLens addresses all three. A daily agent fuses every stream into an
interpretable per-worker report with metrics and flags; a monthly analytics
layer answers the question managers actually ask --- \emph{which conditions
move output and defects?} --- with correlations and threshold-split
findings computed from the data itself. Two constraints shape the design:
decision logic must be deterministic and auditable (an LLM may phrase a
report but never decides anything in it), and the system must respect
workers (aggregate-first analytics, consent, no disciplinary use).

\section{System overview and architecture}
ShiftLens is a Python~3.12 application with a deliberately thin dependency
footprint: FastAPI and uvicorn are the only third-party packages, and
Pearson correlation is computed with the standard library --- no
pandas/numpy. The moving parts:

\begin{itemize}
\item \textbf{Daily ingestion.} One JSON file per day holds the shared
environment record (noise, temperature, PM2.5) and, per worker, wearable
sensor summaries, object-detection events (PPE compliance, phone and idle
events, safety-zone violations), a four-field micro-survey (fatigue, mood,
perceived productivity, free-text blockers), and the day's output (units,
defects, tasks completed).
\item \textbf{Deterministic agent core.} Pure functions compute three
indices in $[0,1]$ --- fatigue, focus, safety --- plus productivity;
threshold rules turn them into flags; a report writer emits one JSON +
Markdown report per worker per day. Every number in a report is produced by
code.
\item \textbf{LLM phrasing layer (optional).} A local LM Studio server
(OpenAI-compatible API) rewrites only the prose summary; the prompt forbids
inventing numbers, and a template fallback produces the summary when the
model is offline. \emph{The LLM phrases; it never decides.}
\item \textbf{Monthly analytics.} Pearson correlations across all
(worker, day) pairs plus threshold-split findings (e.g.\ mean defects on
days above vs.\ below 85~dB), emitted as \texttt{monthly.json} with a
human-readable insight per correlation.
\item \textbf{Dashboard.} A FastAPI single-page app --- vanilla JS with
inline SVG charts and no CDN, so it works fully offline --- showing
correlation scatters with trend lines, team daily units vs.\ fatigue,
per-worker bars, finding cards, and a daily-report browser.
\item \textbf{CLI.} \texttt{simulate} $\to$ \texttt{report} $\to$
\texttt{monthly} $\to$ \texttt{serve}, or \texttt{demo} for the whole
pipeline with a printed findings summary.
\end{itemize}

\section{Data model and simulation}
The dataset covers 8 workers (roles: assembler, packer, inspector, forklift
operator; per-worker \texttt{base\_skill} in $[0.85, 1.20]$ that scales
output) over 30 days starting 2026-09-01, i.e.\ $8 \times 30 = 240$
(worker, day) records. Static worker metadata (id, name, role, shift, base
skill) lives in \texttt{workers.json}; Table~\ref{tab:schema} summarizes
the daily record.

\begin{table}[htbp]
\centering\small
\caption{Daily record (\texttt{data/daily/YYYY-MM-DD.json}): one shared
environment block plus one entry per worker per day.}
\label{tab:schema}
\begin{tabular}{@{}p{2.1cm}p{8.8cm}p{4.7cm}@{}}
\toprule
Group & Fields & Source \\
\midrule
\texttt{environment} & \texttt{noise\_db}, \texttt{temperature\_c},
  \texttt{air\_quality\_pm25} & environmental sensors (shared, per day)\\
\texttt{sensors} & \texttt{hr\_avg}, \texttt{hr\_max}, \texttt{steps},
  \texttt{station\_time\_min} & wearable stream\\
\texttt{detections} & \texttt{ppe\_ok\_pct}, \texttt{phone\_events},
  \texttt{idle\_events}, \texttt{safety\_zone\_violations} &
  object-detection events\\
\texttt{survey} & \texttt{fatigue}, \texttt{mood},
  \texttt{perceived\_productivity}, \texttt{blockers} &
  daily micro-survey (1--5 + free text)\\
\texttt{output} & \texttt{units}, \texttt{defects},
  \texttt{tasks\_completed} & shift output record\\
\bottomrule
\end{tabular}
\end{table}

For the demo, real feeds are replaced by a seeded simulator
(\texttt{random.Random(42)}). The simulator plants a known causal
structure that the analytics must rediscover --- findings are computed
from data, never hardcoded:

\begin{itemize}
\item noisy days (\texttt{noise\_db} $> 85$~dB) push the defect rate up;
\item higher self-reported fatigue pushes units down and defects up;
\item hot days (\texttt{temperature\_c} $> 28\,^\circ$C) reduce PPE
compliance;
\item phone events --- more likely on low-mood days --- reduce units;
\item per-worker base skill scales units; small daily noise everywhere.
\end{itemize}

The planted structure plays the role of ground truth: the monthly layer is
only correct if it recovers these effects from the raw records alone.

\section{Agent design: metrics, flags, LLM polish}
\textbf{Metrics.} All metrics are pure functions of one worker-day record.
The three composite indices blend survey and sensor evidence into $[0,1]$
scores, with $\operatorname{clamp}(x,0,1)=\min(\max(x,0),1)$:
\begin{align}
\texttt{fatigue\_index} &=
  0.6\cdot\frac{\texttt{fatigue}}{5} \;+\;
  0.4\cdot\frac{\min(\texttt{hr\_avg}/100,\;1.2)}{1.2}
  \label{eq:fatigue}\\[2pt]
\texttt{focus\_score} &=
  \operatorname{clamp}\bigl(1 - 0.08\cdot\texttt{phone\_events}
  - 0.03\cdot\texttt{idle\_events},\;0,\;1\bigr)
  \label{eq:focus}\\[2pt]
\texttt{safety\_score} &=
  \operatorname{clamp}\bigl(\texttt{ppe\_ok\_pct}/100
  - 0.25\cdot\texttt{safety\_zone\_violations},\;0,\;1\bigr)
  \label{eq:safety}
\end{align}
Productivity is the raw \texttt{units} count, normalized per worker against
that worker's own 30-day mean before any cross-worker comparison, so a
skilled assembler and a novice are each judged against their own baseline.

\textbf{Flags.} Flags are deterministic threshold rules over the metrics
(Table~\ref{tab:flags}); they are the only mechanism by which a report
raises a hand.

\begin{table}[htbp]
\centering\small
\caption{Deterministic flag rules evaluated on every worker-day.}
\label{tab:flags}
\begin{tabular}{@{}llp{6.9cm}@{}}
\toprule
Flag & Rule & What drives it \\
\midrule
\texttt{high\_fatigue} & $\texttt{fatigue\_index} \ge 0.6$ &
  high survey fatigue and/or elevated heart rate\\
\texttt{low\_focus} & $\texttt{focus\_score} < 0.7$ &
  phone and idle events beyond tolerance\\
\texttt{safety\_risk} & $\texttt{safety\_score} < 0.85$ &
  PPE non-compliance and/or safety-zone violations\\
\texttt{defect\_spike} & $\texttt{defects} \ge 2\times$ worker mean &
  quality outlier against the worker's own baseline\\
\bottomrule
\end{tabular}
\end{table}

\textbf{Report and LLM polish.} The daily report
(\texttt{out/daily\_reports/}$\langle$\textit{date}$\rangle$/%
$\langle$\textit{worker}$\rangle$\texttt{.json} + \texttt{.md}) bundles the
metrics, the triggered flags, and a prose summary. The summary
\emph{always} lists the underlying numbers; the optional LLM (LM Studio,
\texttt{SHIFTLENS\_LMS\_HOST} default \url{http://localhost:1234}, model
\texttt{SHIFTLENS\_LMS\_MODEL} default \texttt{bonsai-27b}) only rephrases
it, under a prompt that forbids inventing numbers. When the model is
unreachable, a template produces the same content. Every factual claim in
a report therefore stays traceable to code, while supervisors still get
natural-language text.

\section{Monthly analytics}
The monthly pass aggregates the 240 worker-day records into
\texttt{out/monthly.json}: a team daily series (units, defects, average
fatigue, noise), per-worker aggregates (total units, defect rate, average
fatigue/focus/safety), correlations, and findings. Correlations are
Pearson's $r$ over all (worker, day) pairs ($n = 240$), computed with the
standard library only:
\[
r \;=\; \frac{\sum_{i=1}^{n}(x_i-\bar{x})(y_i-\bar{y})}
{\sqrt{\sum_{i=1}^{n}(x_i-\bar{x})^2}\;\sqrt{\sum_{i=1}^{n}(y_i-\bar{y})^2}}
\]
Seven predeclared pairs are evaluated: fatigue $\leftrightarrow$ units,
fatigue $\leftrightarrow$ defects, noise $\leftrightarrow$ defects,
temperature $\leftrightarrow$ PPE compliance,
phone events $\leftrightarrow$ units, mood $\leftrightarrow$ units, and
focus score $\leftrightarrow$ units. Each result ships with $x$, $y$, $r$,
$n$, and a generated insight string.

Findings are produced by threshold splits rather than hardcoded
statements: the code bins days (e.g.\ $>85$ vs.\ $\le 85$~dB), computes the
group means and the percent delta, and emits a sentence. In the seeded
reference run described in the design contract, the analytics recover the
planted structure: the seeded reference run shows
fatigue$\leftrightarrow$units at $r = -0.61$ (\emph{``Higher fatigue
strongly tracks lower output''}) and findings such as \emph{``Days above
85~dB average 122\% more defects per worker than quieter days (5.8 vs
2.6)''}.

\section{Evaluation and discussion}
\textbf{Validation.} An end-to-end smoke test runs the demo pipeline into a
temporary data directory and asserts that \texttt{monthly.json} contains
the seven correlations with $|r|>0$ for the planted signals, that daily
reports exist for every worker, and that the dashboard module imports. More
importantly, the simulator's planted causal structure provides real ground
truth: the analytics are judged by whether they rediscover the
noise--defect, fatigue--units, temperature--PPE, and phone--units effects,
in the right direction, with none of them hardcoded.

\textbf{What the dashboard reveals.} On one page a supervisor sees which
conditions correlate with output and defects (scatter plots with trend
lines), how team units and average fatigue move together across the month,
per-worker bars, finding cards in plain language, and a browser that drills
from a finding into any worker's daily report (pick a date, pick a worker,
read the rendered Markdown). The intent is to move the conversation from
end-of-month surprise to same-week action: fix the noisy line, rotate
fatigued workers, double down on PPE compliance on hot days.

\textbf{Limitations.} The demo runs on simulated data; $30$ days $\times$
$8$ workers is a small sample for correlation; Pearson's $r$ captures only
linear association, and the observed correlations are observational, not
causal; the three indices are heuristic composites with hand-set weights;
per-worker normalization hides absolute differences between workers.

\textbf{Privacy and ethics.} The monthly layer is aggregate-first:
management-facing analytics default to team-level findings. Wearable and
survey data require worker consent, and daily reports are a coaching and
self-knowledge tool --- explicitly not an input to disciplinary action. All
processing is local (standard library plus an optional local LLM); no data
leaves the machine, and the dashboard works fully offline.

\section{Conclusion}
ShiftLens shows that a useful agentic system does not need an LLM in the
decision path. A deterministic core turns fragmented daily streams into
interpretable per-worker reports; a monthly analytics layer reconnects
wellbeing and working conditions to output and quality; an optional local
LLM adds fluent prose without ever inventing a fact. Validated against
planted ground truth, the pipeline recovers the effects it should --- and
every number on the dashboard is traceable to code. Next steps are real
sensor connectors, longer observation windows, and richer (still auditable)
models beyond linear correlation.

\end{document}
